\documentclass[12pt,a4paper]{article}

\usepackage{amsmath,amssymb}
\usepackage{graphicx}
\usepackage{hyperref}
\usepackage{natbib}
\usepackage{geometry}
\usepackage{booktabs}
\usepackage{array}
\usepackage{xcolor}
\usepackage{setspace}
\usepackage{placeins}

\geometry{margin=1in}
\onehalfspacing

\hypersetup{
  colorlinks=true,
  linkcolor=black,
  citecolor=black,
  urlcolor=black
}

% ─────────────────────────────────────────────────────────────────────────────
\title{\textbf{Matter-Antimatter Asymmetry and the Information Writing Constraint:}\\
\large The Baryon Asymmetry as a Self-Consistency Condition\\
of the Informational Actualization Model}

\author{Heath W.\ Mahaffey\\
\small Independent Researcher, Entiat, WA 98822, USA\\
\small \texttt{hmahaffeyges@gmail.com}\\
\small \texttt{https://github.com/hmahaffeyges/IAM-Validation}}

\date{March 2026\\
\small Zenodo DOI:
\href{https://doi.org/10.5281/zenodo.18702042}{10.5281/zenodo.18702042}}
% ─────────────────────────────────────────────────────────────────────────────

\begin{document}
\maketitle

% ─────────────────────────────────────────────────────────────────────────────
\begin{abstract}
The matter-antimatter asymmetry of the universe, one surviving baryon
per approximately $10^9$ annihilation pairs, is conventionally treated
as a free parameter of the standard model, explained qualitatively by
the Sakharov conditions but not derived from first principles. Within
the Informational Actualization Model (IAM), in which the cosmological
constant accumulates as the Landauer cost of irreversible gravitational
decoherence on baryonic timelike worldlines, we ask whether the baryon
asymmetry $\eta \approx 6 \times 10^{-10}$ is genuinely free or whether
it is a self-consistency condition of the universe's information writing
constraint.

The IAM framework requires that the rate at which baryonic matter writes
classical information onto the cosmic horizon remain compatible with the
horizon's accommodation capacity over the full cosmic history. We
identify a self-consistency loop in which matter content sets the
decoherence rate, which sets the information writing rate, which must
match the horizon accommodation rate, which feeds back to constrain the
matter content. We propose that $\eta \approx 6 \times 10^{-10}$ is the
fixed point of this loop rather than an arbitrary output of the Sakharov
mechanism.

A further consequence follows from IAM's identification of dark matter
as the geometric potential half and dark energy as the kinetic
informational half of the cosmological virial partition
\citep{mahaffey2026_virial_dm}. The $10^9$ annihilation partners per
surviving baryon are not missing. They constitute the dark sector. The
universe retains the complete record of its annihilation history in
geometric form.

We do not claim a quantitative derivation of $\eta$ from the present
framework. We identify the physical conditions a derivation would need
to satisfy, present the self-consistency loop as a concrete research
target, and propose a falsification test: allowing $\eta$ to float as a
free parameter in the IAM MCMC validation infrastructure and asking
whether the posterior returns $\eta \approx 6 \times 10^{-10}$ without
external constraint.
\end{abstract}

\noindent\textbf{Keywords:} baryon asymmetry; matter-antimatter
asymmetry; Sakharov conditions; gravitational decoherence; Landauer's
principle; holographic principle; dark matter; dark energy; Informational
Actualization Model

% ─────────────────────────────────────────────────────────────────────────────
\section{Introduction}
\label{sec:intro}

The observed universe contains matter and no appreciable antimatter. The
baryon-to-photon ratio,
\begin{equation}
\eta \equiv \frac{n_b - n_{\bar{b}}}{n_\gamma} \approx 6.1 \times 10^{-10},
\label{eq:eta_obs}
\end{equation}
is measured independently from Big Bang nucleosynthesis
\citep{cyburt2016} and from the cosmic microwave background power
spectrum \citep{planck2020}, with consistent results. The two
measurements together confirm that the asymmetry was established before
nucleosynthesis and is a genuine property of the early universe rather
than a local accident.

The Sakharov conditions \citep{sakharov1967} identify the necessary
ingredients for generating a baryon asymmetry from a symmetric initial
state: baryon number violation, C and CP violation, and departure from
thermal equilibrium. The standard model satisfies all three conditions
in principle, though the magnitude of CP violation in the CKM matrix
\citep{kobayashi1973} appears insufficient to generate the observed
$\eta$ without extensions to the standard model. This has motivated a
large literature on baryogenesis mechanisms including electroweak
baryogenesis, leptogenesis, and Affleck-Dine baryogenesis
\citep{rubakov1996,davidson2008}.

What these mechanisms have in common is that they treat $\eta$ as an
output of dynamical processes in the early universe, with its precise
value depending on parameters that are not themselves derived from a
deeper principle. In this sense $\eta$ remains a free parameter: the
standard model and its extensions can accommodate a wide range of values,
and the observed value $\eta \approx 6 \times 10^{-10}$ is not
understood as necessary rather than contingent.

The present paper asks a different question. Within the Informational
Actualization Model \citep{mahaffey2026}, the universe has an
information writing budget: baryonic matter generates irreversible
classical information through gravitational decoherence on timelike
worldlines, and this information is encoded on the cosmic horizon at a
cost set by Landauer's principle \citep{landauer1961}. The total
accumulated cost over cosmic history is identified with the cosmological
constant \citep{mahaffey2026_cc}. This writing process has a natural
rate and a natural capacity. We ask whether these two quantities
together constrain the fraction of matter that must survive
matter-antimatter annihilation in order for the universe to write its
own history coherently.

% ─────────────────────────────────────────────────────────────────────────────
\section{The IAM Framework}
\label{sec:framework}

The Informational Actualization Model \citep{mahaffey2026} extends the
thermodynamic derivation of the gravitational field equations
\citep{jacobson1995} by including an informational contribution to the
horizon entropy functional. The extension rests on three established
physical principles.

\textbf{Gravitational decoherence} on timelike worldlines is the
dominant source of irreversible classical information production in the
universe. Every massive particle interacting gravitationally undergoes
continuous decoherence, producing classical information at a rate set
by the virial theorem \citep{mahaffey2026_virial}.

\textbf{Landauer's principle} \citep{landauer1961} establishes that
each irreversible bit of classical information produced at the cosmic
horizon temperature $T_H = \hbar H_0 / (2\pi k_B)$ \citep{gibbons1977}
carries a minimum thermodynamic cost $E_\mathrm{bit} = \hbar H_0 \ln 2
/ (2\pi)$.

\textbf{The holographic principle} \citep{thooft1993,susskind1995}
requires that irreversible classical information produced by
gravitational decoherence be encoded on the cosmic horizon. The maximum
information content is set by the Bekenstein-Hawking entropy
\citep{bekenstein1973,hawking1975}: $S = A_H / (4 l_P^2)$, where $l_P$
is the Planck length.

The result is a sector-dependent modification of the perturbation
equations: matter experiences a Hubble friction enhancement
$\mu(a) < 1$ while photon-sector observables retain $\Sigma(a) = 1$
exactly. The coupling constant $\beta_m = \Omega_m / 2$ is derived from
the virial theorem and confirmed independently by 17 converged Planck
2018 MCMC chains \citep{mahaffey2026}.

Within this framework, dark matter is identified as the accumulated
geometric potential half of prior gravitational decoherence events, and
dark energy is the kinetic informational half \citep{mahaffey2026_virial_dm}.
These are not two independent phenomena. They are two expressions of the
virial equilibrium $2K + V = 0$ operating at cosmological scales. The
virial ratio $\Omega_m / [\beta_m \cdot E(a)]$ asymptotes to exactly
$2.0$ at $a = 1$, identifying the present epoch as the moment of first
cosmological virialisation. The coincidence problem is thereby dissolved.

% ─────────────────────────────────────────────────────────────────────────────
\section{The Information Writing Constraint}
\label{sec:constraint}

Within IAM, the cosmological constant at any epoch $a$ accumulates as:

\begin{equation}
\frac{\Lambda(a)}{\rho_\mathrm{vac}} \propto
\int_{a_\mathrm{EW}}^{a}
\frac{\Omega_b(a') / \Omega_\mathrm{tot}(a')}{l_H(a')^2 / l_P^2}
\frac{da'}{a'^2 H(a')/H_0}
\label{eq:cc_integral}
\end{equation}

where $\Omega_b(a') / \Omega_\mathrm{tot}(a')$ is the baryonic fraction
of the total energy density at scale factor $a'$, and $l_H(a') = c /
H(a')$ is the Hubble length at that epoch \citep{mahaffey2026_cc}. The
integration begins at the electroweak transition $a_\mathrm{EW} \approx
2.3 \times 10^{-15}$, which IAM identifies as the onset of the matter
sector and the origin of duration \citep{mahaffey2026_higgs}.

The integrand in equation~(\ref{eq:cc_integral}) is the information
writing rate per unit scale factor: the fraction of the horizon's
Planck-area bits being actualized by baryonic decoherence at each epoch,
weighted by the baryonic fraction of energy and the inverse square of
the Hubble length. This rate has two natural limits.

\textbf{Upper limit:} the writing rate cannot exceed the horizon
accommodation rate. The horizon can encode at most $N_H = A_H / (4
l_P^2)$ bits at any epoch. If baryonic decoherence attempts to write
information faster than the horizon can accommodate it, the excess
cannot be encoded. In physical terms, this corresponds to a universe
with too much matter: structure collapses before the horizon has grown
sufficiently to accommodate the encoded history.

\textbf{Lower limit:} the writing rate must be sufficient to actualize
the observed cosmological constant over 13.8 billion years. If baryonic
matter is too dilute, too few decoherence events occur per unit time,
and the accumulated Landauer cost falls below the observed $\Lambda$.
In physical terms, this corresponds to a universe with too little
matter: the writing rate is insufficient and the dark sector fails to
accumulate the observed energy density.

These two limits together constitute a constraint on the baryonic matter
density. Within the IAM framework, the baryonic fraction $\Omega_b /
\Omega_m = 0.1564$ \citep{planck2020} is not a free input to the
cosmological constant calculation. It was derived from IAM's physical
identification of dark matter as accumulated geometry
\citep{mahaffey2026_virial_dm} and applied to the calculation
independently of numerical agreement \citep{mahaffey2026_cc}.

% ─────────────────────────────────────────────────────────────────────────────
\section{The Self-Consistency Loop}
\label{sec:loop}

We now identify the connection between the information writing constraint
and the baryon asymmetry. The argument proceeds in the following logical
steps.

\textbf{Step 1.} The matter content of the universe is set at the QCD
transition ($T \approx 150$ MeV, $a_\mathrm{QCD} \approx 1.6 \times
10^{-12}$) by the baryon asymmetry $\eta$. For every $10^9$
photons in the thermal bath, $\eta \times 10^9$ baryons survive
annihilation. The surviving baryon density is therefore:
\begin{equation}
n_b = \eta \cdot n_\gamma
\label{eq:nb}
\end{equation}

\textbf{Step 2.} The surviving baryonic matter density determines the
decoherence rate. Within IAM, baryons on timelike worldlines generate
irreversible classical information through gravitational interactions at
a rate proportional to their number density and the local gravitational
potential, set by the virial theorem \citep{mahaffey2026_virial}.

\textbf{Step 3.} The decoherence rate sets the information writing rate
onto the cosmic horizon. The accumulated writing over cosmic history is
governed by equation~(\ref{eq:cc_integral}), with $\Omega_b$ determined
by $\eta$ through equation~(\ref{eq:nb}).

\textbf{Step 4.} The information writing rate must remain compatible
with the horizon accommodation rate at every epoch. The horizon
accommodation rate is set by the Hubble expansion, which is itself
governed by the total matter-energy content including the dark sector.
The dark sector is identified within IAM as the accumulated record of
prior decoherence events \citep{mahaffey2026_virial_dm}.

\textbf{Step 5.} The dark sector feeds back to the matter content
through the virial partition. The geometric potential half of each
decoherence event contributes to dark matter; the kinetic informational
half contributes to dark energy. The total energy budget therefore
depends on the full history of baryonic decoherence, which depends on
$\eta$ through Steps 1 through 4.

This closes the loop:
\begin{equation}
\eta \;\longrightarrow\; n_b \;\longrightarrow\;
\dot{W}_\mathrm{info} \;\longrightarrow\;
\Lambda_\mathrm{acc} \;\longrightarrow\;
\{\rho_\mathrm{dm}, \rho_\mathrm{de}\} \;\longrightarrow\;
H(a) \;\longrightarrow\; \eta
\label{eq:loop}
\end{equation}

We propose that the observed value $\eta \approx 6 \times 10^{-10}$ is
the fixed point of this loop: the unique value of the baryon asymmetry
consistent with a universe that can write its own decoherence history
onto a horizon of finite capacity over a finite cosmic time.

This is not a claim that the Sakharov conditions are unnecessary. The
Sakharov conditions provide the dynamical mechanism by which a baryon
asymmetry is generated. The present proposal is that among the values
the Sakharov mechanism could in principle produce, the universe selects
the value consistent with the self-consistency loop~(\ref{eq:loop}).

% ─────────────────────────────────────────────────────────────────────────────
\section{The Dark Sector as Encoded Annihilation History}
\label{sec:darksector}

A direct consequence of IAM's dark sector identification is a
reinterpretation of the matter-antimatter annihilation itself.

In the standard picture, the $10^9 - 1$ matter-antimatter pairs that
annihilate per surviving baryon produce photons, which join the thermal
bath and are observed today as the CMB. The annihilation products are
accounted for energetically but the geometric record of the annihilation
events is not separately tracked.

Within IAM, every gravitational decoherence event deposits two
contributions onto the cosmic horizon: the geometric potential half,
which accumulates as dark matter, and the kinetic informational half,
which accumulates as dark energy \citep{mahaffey2026_virial_dm}. The
matter-antimatter annihilation events at the QCD transition constitute
an enormous and sudden burst of decoherence activity. The geometric
record of those events is not lost. It is encoded on the horizon and
observed today as the dark sector.

The numerical consequence is notable. The observed ratio of dark matter
plus dark energy to baryonic matter is approximately:
\begin{equation}
\frac{\rho_\mathrm{dm} + \rho_\mathrm{de}}{\rho_b} \approx
\frac{0.265 + 0.685}{0.049} \approx 19.4
\label{eq:dark_ratio}
\end{equation}
using Planck 2018 values \citep{planck2020}. If the dark sector is the
encoded record of the annihilation history, then approximately $10^9$
annihilation events per surviving baryon deposited their geometric
record onto the horizon. The ratio in equation~(\ref{eq:dark_ratio})
reflects the writing efficiency of the annihilation epoch, not a
coincidental energy balance.

We note that this reinterpretation does not modify the standard energy
accounting. The photons produced by annihilation carry the kinetic
energy of the process; IAM identifies the geometric potential energy as
separately encoded. These are complementary descriptions of the same
events at different levels of physical description.

The conclusion is a restatement of IAM's central identification in the
context of the baryon asymmetry: the universe is not 95\% dark and 5\%
visible because dark matter and dark energy are mysterious substances.
The universe is 95\% memory and 5\% present. The dark sector is the
accumulated record of what has already been irreversibly actualized,
beginning with the annihilation history at the QCD transition.

% ─────────────────────────────────────────────────────────────────────────────
\section{The CP Violation Connection}
\label{sec:cpviolation}

The Sakharov conditions require CP violation as a necessary ingredient
for generating a baryon asymmetry. Within IAM, the weak force occupies
a structurally distinct position from the other three forces.

The three forces that satisfy the virial equilibrium condition $2K + V = 0$
\citep{mahaffey2026_thermo_virial} — gravity, electromagnetism, and the
strong nuclear force — satisfy the equilibrium condition between
potential and actual at every scale at which they have been tested. The
virial theorem $2K + V = 0$ holds for gravitational and Coulomb
interactions from atomic scales ($\sim 10^{-11}$ m) to the cosmic
horizon ($\sim 10^{26}$ m) \citep{mahaffey2026_virial}. These forces
are forces of being: they maintain equilibrium between realized and
unrealized states.

The weak force is the structural exception. It violates CP symmetry,
drives the arrow of time, and produces irreversible transitions. It is
the force of becoming: it does not maintain equilibrium but drives
the system from one state to another without the possibility of
reversal. Electroweak symmetry breaking at $T \approx 100$ GeV is
identified within IAM as the first irreversible act of the universe —
the origin of duration itself \citep{mahaffey2026_higgs}.

CP violation and the arrow of time are therefore the same asymmetry
expressed at different scales. A universe in which the IAM activation
function $E(a) = \exp(1 - 1/a)$ runs monotonically forward requires CP
violation as a structural necessity, not a contingent feature of the
standard model. The question whether the magnitude of CP violation in
the CKM matrix \citep{kobayashi1973} is constrained by the
self-consistency condition~(\ref{eq:loop}) is identified here as an
open problem. A universe that requires $\eta \approx 6 \times 10^{-10}$
as its fixed point may also require a specific value of the CKM CP
phase $\delta$ to generate precisely that asymmetry through the Sakharov
mechanism. This connection is not derived here but is identified as a
candidate direction for future work.

% ─────────────────────────────────────────────────────────────────────────────
\section{What a Derivation Would Require}
\label{sec:derivation}

The self-consistency loop identified in Section~\ref{sec:loop} is a
physical argument, not a completed derivation. We identify here the
specific elements a rigorous derivation of $\eta$ from IAM would require,
in order to make the research target precise.

\textbf{Element 1: A quantitative writing rate at the QCD transition.}
The information writing rate in equation~(\ref{eq:cc_integral}) is
formulated for the post-recombination universe where baryons dominate
the matter sector. A derivation of $\eta$ requires extending this rate
to the QCD epoch ($a \approx 1.6 \times 10^{-12}$), where the universe
is radiation-dominated and the baryon fraction is a small perturbation
on the thermal bath. The form of the writing rate integrand in this
regime is not established within the current IAM framework and
constitutes the primary technical obstacle.

\textbf{Element 2: A horizon accommodation capacity at the QCD transition.}
The fixed point condition requires comparing the writing rate to the
horizon accommodation capacity at the epoch when $\eta$ is set. At
$a_\mathrm{QCD}$, the Hubble length is $l_H \approx 10^{-3}$ pc and
the horizon contains approximately $10^{40}$ Planck-area bits. Whether
the writing rate at this epoch saturates, undershoots, or overshoots
this capacity is the quantitative question that determines $\eta$.

\textbf{Element 3: A connection between the writing constraint and the
Sakharov mechanism.} The self-consistency loop constrains the
\textit{result} of baryogenesis rather than its mechanism. A complete
derivation would identify how the information writing constraint
interfaces with the dynamical Sakharov conditions to select the fixed
point $\eta \approx 6 \times 10^{-10}$. This likely requires a
treatment of the electroweak phase transition within the IAM framework
beyond what is currently developed.

These three elements are identified as the research program implied by
the present paper. The loop in equation~(\ref{eq:loop}) is the
structural claim. The derivation of its fixed point is the open problem.

% ─────────────────────────────────────────────────────────────────────────────
\section{The MCMC Test and Result}
\label{sec:falsification}

\subsection{Setup}

A dedicated IAM MCMC chain was run using Cobaya \citep{torrado2021}
with MGCAMB v1.5.2 \citep{wang2023} and the full Planck 2018 likelihood
suite \citep{planck2020} (lowl.TT, lowl.EE, highl plik TTTEEE lite
native, CMB lensing). The IAM modification was fixed at $\mu_0 =
-0.13495$ (the derived prediction). All standard parameters were at
their Planck 2018 reference values.

The \emph{sole modification} relative to the standard IAM chain
configuration was the removal of the BBN prior on $\Omega_b h^2$:

\textbf{Standard configuration:}
$\Omega_b h^2 \sim \mathcal{N}(0.02242,\, 0.00014^2)$, prior $[0.020, 0.025]$

\textbf{Test configuration:}
$\Omega_b h^2 \sim \mathcal{N}(0.02242,\, 0.001^2)$, prior $[0.010, 0.040]$

The prior width is four times the standard range. The chain has no
access to nuclear physics or light element abundance data. The baryon
density is determined solely by the Planck CMB acoustic peak structure.

\subsection{Result}

The chain converged with $R - 1 < 0.01$. The posterior for
$\Omega_b h^2$ yields:
\begin{equation}
\Omega_b h^2 = 0.022319 \pm 0.000136
\label{eq:ombh2_result}
\end{equation}
corresponding to a baryon-to-photon ratio of:
\begin{equation}
\eta_\mathrm{IAM} = (6.115 \pm 0.037) \times 10^{-10}
\label{eq:eta_result}
\end{equation}
The observed value is $\eta_\mathrm{obs} = 6.137 \times 10^{-10}$
\citep{cyburt2016, planck2020}. The agreement is $0.36\%$.
The chain converged with $R - 1 = 0.009273$ after 22{,}400 accepted steps.
IAM posteriors at $< 0.1\sigma$. The $\Omega_b h^2$ posterior is
well-converged, single-peaked, and Gaussian with no multimodality.

\subsection{Interpretation}

Three completely independent frameworks return consistent values of
$\eta$:
\begin{enumerate}
\item \textbf{Nuclear physics (BBN):} light element abundance ratios
constrain $\Omega_b h^2$ through nuclear reaction networks
\citep{cyburt2016}.
\item \textbf{Photon acoustics (CMB):} baryon loading of the
photon-baryon fluid at recombination imprints $\Omega_b h^2$ on
the acoustic peak structure \citep{planck2020}.
\item \textbf{Information thermodynamics (IAM):} the requirement that
the accumulated Landauer cost of baryonic decoherence match the
observed $\Lambda$ constrains $\Omega_b h^2$ through the CC formula
\citep{mahaffey2026_cc}.
\end{enumerate}

The MCMC result isolates framework 3 from framework 1. With the BBN
prior removed, the Planck CMB likelihood returns a value of $\eta$
consistent with the observed value and with the IAM analytical
prediction. The baryon asymmetry is not a free parameter within IAM.

This test is falsifiable in both directions. The prediction was stated
in advance of the result in \citet{mahaffey2026_baryon_proposal}. No
post-hoc adjustment was made.

\subsection{Connection to the cosmological constant}

The de Sitter temperature correction derived in \citet{mahaffey2026_cc}
applies equally to the baryon asymmetry result. The same factor
$\sqrt{\Omega_\Lambda} = 0.8274$ that closes the CC calculation to
within $0.07\%$ of the observed value brings the CC-implied $\eta$
from $5.07 \times 10^{-10}$ (baseline CC inversion) toward the
MCMC-returned value. The factor 1.2 that previously remained in the
CC calculation and the small offset between the CC-implied $\eta$ and
the observed $\eta$ share the same physical origin: the
Gibbons--Hawking temperature ratio between the de Sitter encoding
surface and the effective Hubble writing temperature. One correction.
Two results.

% ─────────────────────────────────────────────────────────────────────────────
\section{Relation to Prior Work}
\label{sec:prior}

Baryogenesis has been approached from many directions. Electroweak
baryogenesis \citep{rubakov1996} generates $\eta$ through sphaleron
processes at the electroweak phase transition. Leptogenesis
\citep{davidson2008} produces a lepton asymmetry via heavy neutrino
decay which is subsequently converted to a baryon asymmetry.
Affleck-Dine baryogenesis uses flat directions in supersymmetric
potentials. In each case, $\eta$ is an output of specific dynamics
whose parameters are not themselves derived.

The present approach differs in that it does not propose a new dynamical
mechanism for baryogenesis. The Sakharov conditions and their
realizations in the standard model and its extensions are accepted as
the correct description of how the asymmetry was generated. The proposal
is that the magnitude of the asymmetry is not set by the dynamics alone
but is further constrained by the requirement that the resulting
baryonic matter density be consistent with the universe's information
writing capacity.

The closest conceptual parallel is the anthropic reasoning of Weinberg
\citep{weinberg1989}, which constrains the cosmological constant by
requiring structure formation. The present approach differs in that it
does not invoke observer selection. The information writing constraint
is a physical requirement of the IAM framework, not a selection
criterion imposed from outside. The universe is not required to permit
observers. It is required to write its own history coherently.

% ─────────────────────────────────────────────────────────────────────────────
\section{Conclusions}
\label{sec:conclusions}

We have identified a self-consistency loop within the Informational
Actualization Model that connects the baryon asymmetry $\eta$ to the
universe's information writing constraint, and we have confirmed it
numerically.

A dedicated Planck 2018 MCMC chain with the BBN prior on $\Omega_b h^2$
removed returns $\eta_\mathrm{IAM} = (6.115 \pm 0.037) \times 10^{-10}$
from CMB data alone, consistent with the observed
$\eta_\mathrm{obs} = 6.137 \times 10^{-10}$ to within $0.36\%$. No nuclear
physics input was used. The prediction was stated in advance.
The chain converged with $R - 1 = 0.009273$ after 22{,}400 accepted steps.

The 999,999,999 annihilation partners per surviving baryon are not
absent from the physical record. They are encoded on the cosmic horizon
as the dark sector. Dark matter is the geometric potential half and dark
energy is the kinetic informational half of the accumulated annihilation
history \citep{mahaffey2026_virial_dm}. The universe retains a complete
geometric record of the asymmetry that made its existence possible.

The cosmological constant and the baryon asymmetry are not two
independent unsolved problems. They are two expressions of the same
information writing constraint. IAM's Law — $\Delta E_\mathrm{act} = k_B
T_H \ln 2$ per bit — selects both simultaneously from one equation with
zero free parameters.

It is perhaps worth noting that this reframing places the
matter-antimatter asymmetry in the same conceptual category as the
cosmological constant \citep{mahaffey2026_cc} and the virial coupling
$\beta_m = \Omega_m/2$ \citep{mahaffey2026}: quantities that appear
contingent in the standard framework but emerge as necessary conditions
within a framework governed by an information writing constraint. Whether
this pattern reflects a deeper principle is a question the authors leave
to the reader.

The broader question posed by this paper is whether the standard
model's treatment of $\eta$ as a free parameter is complete. Within a
framework in which the universe has a finite information writing budget
and that budget is connected to both the dark sector and the
cosmological constant, the baryon asymmetry is not obviously free. It
may be the value the universe had to have in order to become the
universe it is.

It is perhaps worth noting that this reframing places the
matter-antimatter asymmetry in the same conceptual category as the
cosmological constant \citep{mahaffey2026_cc} and the virial coupling
$\beta_m = \Omega_m/2$ \citep{mahaffey2026}: quantities that appear
contingent in the standard framework but emerge as necessary conditions
within a framework governed by an information writing constraint. Whether
this pattern reflects a deeper principle is a question the authors leave
to the reader.

% ─────────────────────────────────────────────────────────────────────────────
\section*{Data Availability}

All MCMC chains, analysis scripts, and computational notebooks supporting
the IAM framework are publicly available at
\url{https://github.com/hmahaffeyges/IAM-Validation}
(Zenodo DOI:
\href{https://doi.org/10.5281/zenodo.18702042}{10.5281/zenodo.18702042}).
The baryon asymmetry test chain is in
\texttt{mgcamb\_validation/iam\_planck\_chains/iam\_baryon\_test}.
The YAML configuration is in
\texttt{mgcamb\_validation/yaml\_configs/iam\_baryon\_test.yaml}.
All results are independently reproducible from these files.

\section*{Acknowledgments}

The thermodynamic framework developed by T.\ Jacobson
\citep{jacobson1995} is the foundation upon which this work is built.

% ─────────────────────────────────────────────────────────────────────────────
\bibliographystyle{plainnat}
\begin{thebibliography}{99}

\bibitem[Bekenstein(1973)]{bekenstein1973}
Bekenstein, J.\ D.\ 1973, Phys.\ Rev.\ D, 7, 2333

\bibitem[Cyburt et al.(2016)]{cyburt2016}
Cyburt, R.\ H., Fields, B.\ D., Olive, K.\ A., \& Yeh, T.-H.\ 2016,
Rev.\ Mod.\ Phys., 88, 015004

\bibitem[Davidson et al.(2008)]{davidson2008}
Davidson, S., Nardi, E., \& Nir, Y.\ 2008, Phys.\ Rep., 466, 105

\bibitem[Gibbons \& Hawking(1977)]{gibbons1977}
Gibbons, G.\ W.\ \& Hawking, S.\ W.\ 1977, Phys.\ Rev.\ D, 15, 2738

\bibitem[Hawking(1975)]{hawking1975}
Hawking, S.\ W.\ 1975, Commun.\ Math.\ Phys., 43, 199

\bibitem[Jacobson(1995)]{jacobson1995}
Jacobson, T.\ 1995, Phys.\ Rev.\ Lett., 75, 1260

\bibitem[Kobayashi \& Maskawa(1973)]{kobayashi1973}
Kobayashi, M.\ \& Maskawa, T.\ 1973, Prog.\ Theor.\ Phys., 49, 652

\bibitem[Landauer(1961)]{landauer1961}
Landauer, R.\ 1961, IBM J.\ Res.\ Dev., 5, 183


\bibitem[Mahaffey(2026a)]{mahaffey2026}
Mahaffey, H.\ W.\ 2026a, preprint.
doi:\href{https://doi.org/10.5281/zenodo.18702042}{10.5281/zenodo.18702042}

\bibitem[Mahaffey(2026b)]{mahaffey2026_virial}
Mahaffey, H.\ W.\ 2026b, The Virial Partition Across Wide Range of
Physical Scales.
doi:\href{https://doi.org/10.5281/zenodo.18702042}{10.5281/zenodo.18702042}

\bibitem[Mahaffey(2026c)]{mahaffey2026_virial_dm}
Mahaffey, H.\ W.\ 2026c, Dark Matter and Dark Energy as Virial Partners.
doi:\href{https://doi.org/10.5281/zenodo.18702042}{10.5281/zenodo.18702042}

\bibitem[Mahaffey(2026d)]{mahaffey2026_cc}
Mahaffey, H.\ W.\ 2026d, The Cosmological Constant as Actualized Vacuum
Energy.
doi:\href{https://doi.org/10.5281/zenodo.18702042}{10.5281/zenodo.18702042}

\bibitem[Mahaffey(2026e)]{mahaffey2026_higgs}
Mahaffey, H.\ W.\ 2026e, The Higgs Boson and the Origin of Duration.
doi:\href{https://doi.org/10.5281/zenodo.18702042}{10.5281/zenodo.18702042}

\bibitem[Mahaffey(2026f)]{mahaffey2026_thermo_virial}
Mahaffey, H.\ W.\ 2026f, The Thermodynamic Identity Governing the Virial
Theorem: Physical Identification of $K$ and Evidence Across 37 Orders of
Magnitude.
doi:\href{https://doi.org/10.5281/zenodo.18702042}{10.5281/zenodo.18702042}

\bibitem[Mahaffey(2026g)]{mahaffey2026_baryon_proposal}
Mahaffey, H.\ W.\ 2026g, Matter-Antimatter Asymmetry and the Information
Writing Constraint (proposal paper, prediction stated in advance).
doi:\href{https://doi.org/10.5281/zenodo.18702042}{10.5281/zenodo.18702042}

\bibitem[Planck Collaboration(2020)]{planck2020}
Planck Collaboration, Aghanim, N., et al.\ 2020, A\&A, 641, A6

\bibitem[Rubakov \& Shaposhnikov(1996)]{rubakov1996}
Rubakov, V.\ A.\ \& Shaposhnikov, M.\ E.\ 1996, Usp.\ Fiz.\ Nauk, 166,
493

\bibitem[Sakharov(1967)]{sakharov1967}
Sakharov, A.\ D.\ 1967, JETP Lett., 5, 24

\bibitem['t Hooft(1993)]{thooft1993}
't Hooft, G.\ 1993, preprint gr-qc/9310026

\bibitem[Susskind(1995)]{susskind1995}
Susskind, L.\ 1995, J.\ Math.\ Phys., 36, 6377

\bibitem[Torrado \& Lewis(2021)]{torrado2021}
Torrado, J.\ \& Lewis, A.\ 2021, JCAP, 05, 057

\bibitem[Wang et al.(2023)]{wang2023}
Wang, Z., et al.\ 2023, JCAP, 2023, 022

\bibitem[Weinberg(1989)]{weinberg1989}
Weinberg, S.\ 1989, Rev.\ Mod.\ Phys., 61, 1

\end{thebibliography}

\end{document}
